\section*{Chapitre \ref{ch:nw:cross-region-crypto-networks} --- Réseaux crypto inter-régions}

\begin{solution}{exo:nw:cross-region-crypto-networks:1}
Face à un plancher aller-retour sur fibre de \qty{91.3}{\milli\second}~: $224.3/91.3 = 2.46$ pour le backbone et $136.5/91.3 = 1.50$ pour le spécialiste.
\end{solution}

\begin{solution}{exo:nw:cross-region-crypto-networks:2}
Le BGP choisit entre les routes selon les politiques des réseaux qui les annoncent (relations commerciales, préférences, longueur du chemin en \omterm{def:nw:where-the-crypto-venues-live:as}{systèmes autonomes}),
et non selon le délai mesuré~; il ne connaît pas du tout la latence d'un chemin.
\end{solution}

\begin{solution}{exo:nw:cross-region-crypto-networks:3}
Francfort $224.3 - 136.5 = 87.8$~ms~; Londres $210.0 - 139.6 = 70.4$~ms~; Stockholm $245.3 - 124.3 = 121.0$~ms d'aller-retour.
\end{solution}

\begin{solution}{exo:nw:cross-region-crypto-networks:4}
Il corrèle leurs délais et leurs pannes~: une congestion ou une panne sur le câble partagé touche les deux, de sorte que le minimum des deux n'est souvent pas meilleur
que chacune des deux, et la \omterm{def:nw:cross-region-crypto-networks:idem}{mise en concurrence des routes} économise peu précisément dans les événements qu'elle devait couvrir.
\end{solution}

\begin{solution}{exo:nw:cross-region-crypto-networks:5}
Sans risque~: les messages dont la répétition n'a aucun effet (annulations, qui échouent sans conséquence une fois l'ordre disparu~; requêtes~; heartbeats), et les ordres qui restent au carnet assez longtemps
pour que chaque copie arrive pendant que la première est encore ouverte. Risqués~: les ordres qui peuvent s'exécuter à leur arrivée (ordres preneurs de liquidité, immediate-or-cancel), dont les
copies tardives sont acceptées comme de nouveaux ordres une fois le premier entièrement exécuté.
\end{solution}

\begin{solution}{exo:nw:cross-region-crypto-networks:6}
Au plus près de la plateforme dont il doit réagir le plus vite aux prix, généralement celle où il cote ou prend, les données de l'autre marché lui étant
acheminées par le chemin le plus rapide~; s'il couvre sur les deux marchés, un moteur dans chaque région, reliés par une liaison rapide. L'aller-retour de l'ordre de \qty{100}{\milli\second} signifie que le
moteur doit se trouver là où se joue la course.
\end{solution}

\begin{solution}{exo:nw:cross-region-crypto-networks:7}
Entre $\rho = 0.46$ (\qty{110}{\micro\second}) et $\rho = 0.48$ (\qty{98}{\micro\second})~: au-delà d'environ 0.48, la troisième route économise moins que
\qty{100}{\micro\second} au 99\up{e}~centile.
\end{solution}

\begin{solution}{exo:nw:cross-region-crypto-networks:8}
Une plateforme qui ne rejette les doublons que tant que le premier ordre est ouvert accepte une copie tardive une fois le premier entièrement exécuté~: pour les ordres preneurs de liquidité, la mise en concurrence sur trois
routes peut signifier trois exécutions. Elle ne constitue une latence gratuite que pour les messages que la plateforme traite de façon idempotente.
\end{solution}

\begin{solution}{pb:nw:cross-region-crypto-networks:1}
\textbf{Partie I.}
\begin{enumerate}
\item \qty{5311}{\kilo\metre}~; planchers aller-retour de \qty{35.4}{\milli\second} dans le vide et \qty{51.8}{\milli\second} sur fibre.
\item 1.26 pour le spécialiste, 1.31 pour le backbone.
\item $65.4 - 51.8 = \qty{13.6}{\milli\second}$ d'aller-retour.
\item Parce que personne ne l'a publié~: c'est l'hypothèse du chapitre pour un chemin choisi par le BGP, plus lent et à queue plus épaisse.
\end{enumerate}
\textbf{Partie II.}
\begin{enumerate}[resume]
\item 1\,623, 771 et \qty{53}{\micro\second}.
\item 2\,141, 861 et \qty{53}{\micro\second}.
\item À 10~USD par microseconde et par mois~: 5\,182, 902 et 0~USD par mois.
\item Elle la réduit aussi, de \qty{536}{\micro\second} à $\rho = 0$ avec deux routes et de 290 pour 0.5, car le backbone bat parfois le
  spécialiste~; mais le gain est le plus grand dans la queue de la distribution.
\end{enumerate}
\textbf{Partie III.}
\begin{enumerate}[resume]
\item Deux~: chaque copie tardive.
\item 1.57, 0.50 et 0.07 en moyenne.
\item Annulations, modifications lorsque la plateforme les traite de façon idempotente, et requêtes~; pas les ordres preneurs de liquidité.
\item Qu'un identifiant d'ordre client soit rejeté pendant toute la durée de la session (ou d'une journée), et non seulement tant que l'ordre est ouvert.
\end{enumerate}
\textbf{Partie IV.}
\begin{enumerate}[resume]
\item \emph{Résultat principal}~: la mise en concurrence de deux routes fait gagner \qty{1623}{\micro\second} au 99\up{e}~centile lorsqu'elles sont indépendantes, 771 pour $\rho = 0.5$ et 53 pour
  0.9~; une troisième en fait gagner 518, 90 et 0 de plus, de sorte qu'à 10~USD par microseconde et par mois, la mise en concurrence cesse d'être rentable au-delà de 5\,182, 902 et 0~USD par mois.
\item La corrélation $\rho$.
\item À partir des copies journalisées~: envoyer des sondes (ou chaque message) sur toutes les routes avec horodatage et calculer la corrélation des délais des routes, avec soin pour
  les queues de distribution.
\item Elle bondit si les deux routes empruntent le chemin de remplacement du câble coupé, ou diminue si une route passe sur un chemin séparé~; il faut la remesurer après chaque événement.
\item Mettre en concurrence le backbone avec un spécialiste \qty{88}{\milli\second} plus rapide n'économise rien, sauf si le spécialiste tombe en panne~; utiliser le backbone comme route de secours,
  et non comme route en concurrence.
\item Tout s'y ajoute~: une instance lente à l'une ou l'autre extrémité gâche la route la plus rapide, donc la \omterm{def:nw:where-the-crypto-venues-live:lottery}{loterie des instances} passe en premier.
\item Dans les lignes inter-régions~: les routes, les prix, la règle de mise en concurrence par type de message et les règles de doublons des plateformes.
\item Lorsque les routes tombent en panne de façon indépendante, la queue de distribution compte pour la stratégie, et la plateforme n'exécutera pas deux fois une copie tardive.
\end{enumerate}
\end{solution}

\begin{solution}{iq:nw:cross-region-crypto-networks:1}
Le chemin du backbone est choisi pour la capacité et la résilience, non pour la latence, et peut faire un long détour~; on ajoute l'équipement et le routage entre les
zones de la région et le backbone. Les chiffres publiés le placent à 2.5 fois le plancher sur fibre, là où un spécialiste se situe à 1.5.
\iqlookfor{\omterm{def:nw:cross-region-crypto-networks:infl}{Allongement du chemin}~; routage selon les politiques~; le plancher comme référence.}
\end{solution}

\begin{solution}{iq:nw:cross-region-crypto-networks:2}
Lorsque l'edge de la stratégie repose sur la vision des prix d'une autre région plus tôt que les concurrents (arbitrage entre plateformes, couverture), lorsque le gain se chiffre en dizaines
de millisecondes comme entre Tokyo et l'Europe, et lorsque le volume échangé le rentabilise~; non lorsque le backbone est déjà proche du plancher, comme pour
Singapour.
\iqlookfor{Valeur de la latence face au coût~; où le backbone est médiocre~; mesure.}
\end{solution}

\begin{solution}{iq:nw:cross-region-crypto-networks:3}
Le minimum des deux délais~: une médiane plus basse et une queue bien plus basse si les routes tombent en panne de façon indépendante, peu si elles sont corrélées. Cela dépend de la
corrélation, de la queue de chaque route et de la possibilité pour le récepteur d'écarter sans risque la seconde copie.
\iqlookfor{Statistiques d'ordre~; corrélation~; idempotence.}
\end{solution}

\begin{solution}{iq:nw:cross-region-crypto-networks:4}
Une \omterm{def:nw:cross-region-crypto-networks:idem}{clé d'idempotence} attachée à chaque copie d'une même requête afin que le récepteur n'agisse qu'une fois. Elle échoue si le récepteur oublie la clé trop tôt (par
exemple, une fois l'ordre entièrement exécuté), si l'émetteur réutilise des identifiants, ou si des copies différentes portent des clés différentes.
\iqlookfor{Mémoire côté récepteur~; durée de vie des clés~; discipline de l'émetteur.}
\end{solution}

\begin{solution}{iq:nw:cross-region-crypto-networks:5}
Cela rend les stratégies entre plateformes locales~: un moteur doit se trouver à côté de la plateforme qu'il négocie, voir l'autre marché avec un dixième de seconde de retard, et ne négocier que
les signaux qui survivent à ce délai~; l'arbitrage rapide entre régions est une course aux routes, gagnée par ceux qui ont la meilleure ligne privée.
\iqlookfor{Placement du moteur~; horizon des signaux~; la valeur des liaisons privées.}
\end{solution}

\begin{solution}{iq:nw:cross-region-crypto-networks:6}
On ne met en concurrence que les messages idempotents~; pour les nouveaux ordres, on envoie une seule fois sur la meilleure route, avec un basculement rapide~; si une plateforme retient les clés pour la session, on met aussi en concurrence les
ordres avec une seule clé~; on ne génère jamais de nouvelle clé pour une copie~; on rapproche chaque accusé de réception de la clé, et on annule sans délai toute
exposition liée à une exécution en double inattendue.
\iqlookfor{Classes de messages~; règle de la plateforme~; rapprochement~; basculement plutôt que mise en concurrence.}
\end{solution}
